\documentclass[10pt]{article}
\usepackage[T1]{fontenc}
\usepackage[utf8]{inputenc}
\usepackage[french]{babel}
\usepackage{amsmath,amssymb,mathrsfs,mathpazo}
\usepackage{geometry}
\geometry{a4paper,margin=2cm}
\usepackage{graphicx,array,tabularx,tabularray,colortbl,multirow,multicol,booktabs,enumitem,pifont,ulem,fancybox,diagbox,cancel}
\usepackage{pstricks,pstricks-add,pst-plot,pst-tree,pst-node,pst-text,pst-func,pst-math,pst-bspline,pst-3dplot,pst-eucl,pst-solides3d}
\let\pscalculate\undefined
\usepackage{pst-bezier}
\usepackage{tikz,pgfplots,tkz-tab}
\usetikzlibrary{arrows,arrows.meta,patterns,calc,positioning,trees,decorations.markings,shapes.geometric,shapes.misc,shapes.arrows,intersections,decorations.pathmorphing,decorations.pathreplacing,backgrounds,shadows,babel}
\pgfplotsset{compat=1.18}
\usepgfplotslibrary{fillbetween}
\usepackage[np]{numprint}
\usepackage{listings,makecell,pas-tableur,twoopt}
\let\barre\undefined
\usepackage[pstricks]{tablvar}
\definecolor{darkgreen}{RGB}{0,100,0}
\providecommand{\red}{\color{red}}
\providecommand{\blue}{\color{blue}}
\providecommand{\green}{\color{green}}
\providecommand{\black}{\color{black}}
\providecommand{\white}{\color{white}}
\providecommand{\gray}{\color{gray}}
\providecommand{\cyan}{\color{cyan}}
\providecommand{\magenta}{\color{magenta}}
\providecommand{\yellow}{\color{yellow}}
\providecommand{\violet}{\color{violet}}
\providecommand{\orange}{\color{orange}}
\providecommand{\euro}{\texteuro}
\providecommand{\EUR}{\texteuro}
\providecommand{\ieme}{\textsuperscript{e}}
\providecommand{\degres}{\ensuremath{^\circ}}
\providecommand{\degree}{\ensuremath{^\circ}}
\providecommand{\vv}[1]{\overrightarrow{#1}}
\providecommand{\Coord}[1]{#1}
\providecommand{\decosix}{}
\providecommand{\starredbullet}{\ensuremath{\bullet}}
\pagestyle{empty}
\setlength{\parindent}{0pt}
\usepackage{adjustbox}\DeclareUnicodeCharacter{2236}{:}\DeclareUnicodeCharacter{2190}{\ensuremath{\leftarrow}}
\providecommand{\R}{\mathbb{R}}
\providecommand{\N}{\mathbb{N}}
\providecommand{\Z}{\mathbb{Z}}
\providecommand{\Q}{\mathbb{Q}}
\providecommand{\D}{\mathbb{D}}
\providecommand{\C}{\mathbb{C}}
\providecommand{\e}{{\mathrm{e}}}
\providecommand{\vect}[1]{\overrightarrow{#1}}
\providecommand{\vectt}[1]{\overrightarrow{\mathrm{#1}}}
\providecommand{\vvt}[1]{\overrightarrow{\mathrm{#1}}}
\providecommand{\np}[2][]{{#2~#1}}
\providecommand{\vv}[1]{{\overrightarrow{#1}}}
\renewcommand{\d}{{\mathrm{d}}}
\providecommand{\degres}{{^\circ}}
\providecommand{\degre}{{^\circ}}
\providecommand{\blue}{\color{blue}}
\providecommand{\red}{\color{red}}
\providecommand{\enskip}{\quad}
\providecommand{\llbracket}{\lbrack\!\lbrack}
\providecommand{\rrbracket}{\rbrack\!\rbrack}
\providecommand{\euro}{{\text{€}}}
\providecommand{\no}{{\text{n°}}}
\providecommand{\interval}[2]{{[#1;#2]}}
\providecommand{\intervFO}[2]{{[#1;#2[}}
\providecommand{\dotfill}{\ldots}

\providecommand{\ds}{\displaystyle}
\providecommand{\textding}[1]{#1}
\providecommand{\Oij}{\(\mathrm{O};\vec{\imath},\vec{\jmath}\)}
\providecommand{\Oijk}{\(\mathrm{O};\vec{\imath},\vec{\jmath},\vec{k}\)}
\providecommand{\Ouv}{\(\mathrm{O};\vec{u},\vec{v}\)}
\providecommand{\intervFF}[2]{\ensuremath{[#1;#2]}}
\renewcommand{\euro}{\ifmmode\text{\texteuro}\else\texteuro\fi}
\DeclareUnicodeCharacter{00B0}{\ifmmode^{\circ}\else\textdegree\fi}
\DeclareUnicodeCharacter{20AC}{\ifmmode\text{\texteuro}\else\texteuro\fi}
% Énoncés : Éducation nationale. Transcription et corrigés : APMEP, Denis Vergès et contributeurs. Découpage et classement : CoopMaths, Sébastien Lozano, Éric Elter et contributeurs. Adaptation de la présentation : Maths Explorer.
% CC BY-SA : https://coopmaths.fr/www/about/
% https://coopmaths.fr/alea/static/e3c/2020/tex/e3c_2020_00_sujet21_4.tex
\begin{document}
\noindent\textit{APMEP / CoopMaths. CC BY-SA. Présentation adaptée par Maths Explorer.}\par\medskip
\medskip

Une entreprise de \np{1000} employés est organisée en 3 services \og A \fg, \og B \fg{} et \og C \fg{} d’effectifs
respectifs $450$, $230$ et $320$ employés. Une enquête effectuée auprès de tous les employés
sur leur temps de parcours quotidien entre leur domicile et l’entreprise a montré que :

\begin{itemize}
\item 40\,\% des employés du service \og A \fg{} résident à moins de 30 minutes de l’entreprise ;
\item 20\,\% des employés du service \og B \fg{} résident à moins de 30 minutes de l’entreprise ;
\item 80\,\% des employés du service \og C \fg{} résident à moins de 30 minutes de l’entreprise.
\end{itemize}

On choisit au hasard un employé de cette entreprise et on considère les évènements
suivants :

\begin{itemize}
\item $A$ : l’employé fait partie du service \og A \fg{} ;
\item $B$ : l’employé fait partie du service \og B \fg{} ;
\item $C$ : l’employé fait partie du service \og C \fg{} ;
\item $T$ : l’employé réside à moins de $30$~minutes de l’entreprise.
\end{itemize}
On rappelle que si $E$ et $F$ sont deux évènements, la probabilité d’un évènement $E$ est notée
$p(E)$ et celle de $E$ sachant $F$ est notée $p_F(E)$.

\medskip

\begin{enumerate}
\item Justifier que $p(A) = 0,45$ puis donner $p_A(T)$.
\item Compléter l’arbre pondéré ci-dessous.
\psset{nodesepA=0pt,nodesepB=3pt,treesep=0.75,labelsep=0.1pt,levelsep=2.75cm}
\pstree[treemode=R]{\TR{}}
{\pstree{\TR{$A$~}\taput{$\np{0.45}$}}
	{
	\TR{$T$}\taput{$\dots$}
	\TR{$\overline{T}$}\tbput{$\dots$}
	}
\pstree{\TR{$B$~}\tbput{$\dots$}}
	{\TR{$T$}\taput{$\dots$}
	\TR{$\overline{T}$}\tbput{$\dots$}
	}
\pstree{\TR{$C$~}\tbput{$\dots$}}
	{\TR{$T$}\taput{$\dots$}
	\TR{$\overline{T}$}\tbput{$\dots$}
	}	
}
\item Déterminer la probabilité que l'employé choisi soit du service \og A \fg et qu’il réside à moins de
30 minutes de son lieu de travail.
\item Montrer que $p(T) = 0,482$.
\item Sachant qu'un employé de l’entreprise réside à moins de 30 minutes de son lieu de travail,
déterminer la probabilité qu'il fasse partie du service \og C \fg. Arrondir à $10^{-3}$ près.

\end{enumerate}

\end{document}
